%
% connection_lease_reaping_no_release_buggy.tex
%
% FIDELITY CONTROL, not the design. Changes only Reap: it still drops the
% lapsed connection from `registered`, but it never releases the scenes
% that connection owned. This is not a hypothetical weakening -- it is the
% literal current behaviour of HubDisplay.drop_connection and
% HubClientRegistry.named_sessions's sweep, neither of which today calls
% into OwnerTracker at all. See docs/connection_lease_reaping.tex, Section
% "Fidelity: reproducing the shipped defect", for the full trace and the
% invariant checks that must return FOUND against this variant (and NOT
% found against the correct spec).
%
% Type-check:  fuzz docs/connection_lease_reaping_no_release_buggy.tex
% Model-check:
%   probcli docs/connection_lease_reaping_no_release_buggy.tex \
%       -model_check -nodead -goal "..." -p DEFAULT_SETSIZE 3
%

\documentclass[a4paper,10pt,fleqn]{article}
\usepackage[margin=1in]{geometry}
\usepackage{fuzz}
\usepackage[colorlinks=true,linkcolor=blue,citecolor=blue,urlcolor=blue]{hyperref}

\begin{document}

\title{lux Connection-Lease Reaping --- Fidelity Control (No Release)}
\author{Buggy variant: Reap never releases the departing connection's scenes}
\date{September 2026}
\maketitle

% ============================================================
\section{Introduction}
% ============================================================

This is the fidelity control for \texttt{docs/connection\_lease\_reaping
.tex}, not the design. It removes only the ownership-release half of
\texttt{Reap}: a lapsed connection still correctly leaves
\texttt{registered} (I1's mechanism is untouched), but every scene it
owned is left exactly as it was --- this models
\texttt{HubClientRegistry.named\_sessions}'s sweep, which drops a session
from its own dict and never calls into \texttt{OwnerTracker}, composed
with \texttt{HubDisplay.drop\_connection}, whose own docstring says the
departing connection's roots ``stay owned by its id.'' Read the main
specification first; this document states only what differs, in
Section~\ref{sec:reap}.

% ============================================================
\section{Basic Types}
% ============================================================

\begin{zed}
  [CONN, IDENT, SCENE]
\end{zed}

Unchanged from the main specification: three connections, two identities,
and three scenes, bounded by ProB's \texttt{DEFAULT\_SETSIZE}.

% ============================================================
\section{Free Types}
% ============================================================

\begin{zed}
  TSTATE ::= tup | tdown
\end{zed}

\begin{zed}
  LSTATE ::= llive | llapsed
\end{zed}

\begin{zed}
  OWNER ::= unowned | conn \ldata CONN \rdata
\end{zed}

Unchanged from the main specification.

% ============================================================
\section{State}
% ============================================================

\begin{schema}{ConnReg}
  registered : \power CONN \\
  identity : CONN \pfun IDENT \\
  transport : CONN \pfun TSTATE \\
  lease : CONN \pfun LSTATE \\
  owner : SCENE \fun OWNER
\where
  registered \subseteq \dom identity \\
  \dom identity = \dom transport \\
  \dom transport = \dom lease
\end{schema}

Unchanged from the main specification --- see there for the full
description of each field.

% ============================================================
\section{Initialization}
% ============================================================

\begin{schema}{Init}
  ConnReg'
\where
  registered' = \emptyset \\
  identity' = \emptyset \\
  transport' = \emptyset \\
  lease' = \emptyset \\
  owner' = (\lambda s : SCENE @ unowned)
\end{schema}

Unchanged from the main specification.

% ============================================================
\section{Operations}
% ============================================================

\subsection{Reap --- BUGGY: no ownership release}
\label{sec:reap}

\textbf{This is the one operation that differs from the design.} The
lapsed connection \texttt{c?} still correctly leaves \texttt{registered}
--- the lease sweep itself is not the defect, and I1's mechanism is
deliberately left intact by this variant, isolating the I2/I3 finding from
an I1 one. But \texttt{owner} is left completely unchanged: no scene
\texttt{c?} owned is released. Compare against the correct \texttt{Reap}
in \texttt{connection\_lease\_reaping.tex}: the correct version clears
every scene owned by \texttt{c?} to \texttt{unowned} in the same step this
one only drops \texttt{c?} from \texttt{registered}.

\begin{schema}{Reap}
  \Delta ConnReg \\
  c? : CONN
\where
  c? \in registered \\
  lease~c? = llapsed \\
  registered' = registered \setminus \{ c? \} \\
  owner' = owner \\
  identity' = identity \\
  transport' = transport \\
  lease' = lease
\end{schema}

\subsection{Connect, Renew, TransportDies, LeaseLapse, Claim --- unchanged}

Identical to the main specification: only \texttt{Reap}'s
ownership-release step is removed.

\begin{schema}{Connect}
  \Delta ConnReg \\
  c? : CONN \\
  i? : IDENT
\where
  registered' = registered \cup \{ c? \} \\
  identity' = identity \oplus \{ c? \mapsto i? \} \\
  transport' = transport \oplus \{ c? \mapsto tup \} \\
  lease' = lease \oplus \{ c? \mapsto llive \} \\
  owner' = owner
\end{schema}

\begin{schema}{Renew}
  \Delta ConnReg \\
  c? : CONN
\where
  c? \in registered \\
  transport~c? = tup \\
  lease' = lease \oplus \{ c? \mapsto llive \} \\
  registered' = registered \\
  identity' = identity \\
  transport' = transport \\
  owner' = owner
\end{schema}

\begin{schema}{TransportDies}
  \Delta ConnReg \\
  c? : CONN
\where
  c? \in registered \\
  transport~c? = tup \\
  transport' = transport \oplus \{ c? \mapsto tdown \} \\
  registered' = registered \\
  identity' = identity \\
  lease' = lease \\
  owner' = owner
\end{schema}

\begin{schema}{LeaseLapse}
  \Delta ConnReg \\
  c? : CONN
\where
  c? \in registered \\
  transport~c? = tdown \\
  lease~c? = llive \\
  lease' = lease \oplus \{ c? \mapsto llapsed \} \\
  registered' = registered \\
  identity' = identity \\
  transport' = transport \\
  owner' = owner
\end{schema}

\begin{schema}{Claim}
  \Delta ConnReg \\
  c? : CONN \\
  s? : SCENE
\where
  c? \in registered \\
  owner~s? = unowned \lor owner~s? = conn(c?) \\
  owner' = owner \oplus \{ s? \mapsto conn(c?) \} \\
  registered' = registered \\
  identity' = identity \\
  transport' = transport \\
  lease' = lease
\end{schema}

% ============================================================
\section{Verification: the goals that must flip}
% ============================================================

Against this variant, I2's negation must be \emph{found} --- the trace in
\texttt{connection\_lease\_reaping.tex}, Section ``Fidelity: reproducing
the shipped defect,'' reaches it directly:
\texttt{Connect}(c1, i1); \texttt{Claim}(c1, s1); \texttt{TransportDies}(c1);
\texttt{LeaseLapse}(c1); \texttt{Reap}(c1).

\begin{verbatim}
  # must be FOUND against this variant (NOT found against the correct spec)
  probcli docs/connection_lease_reaping_no_release_buggy.tex \
    -model_check -nodead \
    -goal "#s.(#c.(s:SCENE & c:CONN & owner(s) = conn(c) &
                   not(c : registered)))" \
    -p DEFAULT_SETSIZE 3
\end{verbatim}

I3's negation must also be \emph{found}, continuing the same trace one
step further with a live reconnect under the same identity --
\texttt{Connect}(c2, i1):

\begin{verbatim}
  # must be FOUND against this variant (NOT found against the correct spec)
  probcli docs/connection_lease_reaping_no_release_buggy.tex \
    -model_check -nodead \
    -goal "#s.(#c.(#c2.(s:SCENE & c:CONN & c2:CONN &
                        c : registered & not(c2 : registered) &
                        c2 : dom(identity) & identity(c) = identity(c2) &
                        not(c = c2) & owner(s) = conn(c2))))" \
    -p DEFAULT_SETSIZE 3
\end{verbatim}

I1 must stay unaffected --- this variant deliberately leaves the lease
sweep's own removal from \texttt{registered} intact, so it proves nothing
about I1 on its own and must not be mistaken for a stuck-registration
counterexample:

\begin{verbatim}
  # must remain FOUND (the chain still completes -- this variant does not
  # touch the lease sweep's own removal from `registered`, only what
  # happens to `owner` afterward)
  probcli docs/connection_lease_reaping_no_release_buggy.tex \
    -model_check -nodead \
    -goal "#c.(c:CONN & c : dom(identity) & not(c : registered) &
               transport(c) = tdown)" \
    -p DEFAULT_SETSIZE 3
\end{verbatim}

\end{document}
